% GENERATED FILE. Edit canonical lessons, then run npm run generate:books.
% canonical-source-hash: fnv1a64:29a5edd9c2351f01
% canonical-lessons: AR-C136-timsah, AR-C136-bauda, AR-C136-dhubaba, AR-C136-nahla, AR-C136-farasha

\chapter{Crocodile, Mosquito, Fly, Bee, Butterfly}
\label{ch:ar-crocodile-mosquito-fly-bee-butterfly}
\hlchaptermodality{136}

\begin{chapteropening}
\textbf{By the end of this chapter:} \emph{I can say a crocodile, a mosquito, a fly, a bee and a butterfly.}

Complete the last lesson of chapter 136: I can say a crocodile, a mosquito, a fly, a bee and a butterfly.
\end{chapteropening}

\section[\texorpdfstring{\ar{تمساح} (timsāḥ)}{تمساح (timsāḥ)}]{\ar{تمساح} (timsāḥ) — a crocodile}

\label{lesson:AR-C136-timsah}



\begin{quote}
Before the new one: say the Arabic for a duck, then the Arabic for a turtle.

\emph{Read it:} each word whole, right to left — \textbf{\ar{سلام}}\ar{،} \textbf{\ar{شكرًا}}\ar{،} \textbf{\ar{شاي}}
\end{quote}

\subsection*{You'll want to know: \ar{تمساح}}
\textbf{\ar{تمساح}} — \emph{timsāḥ} — \textquotedblleft{}a crocodile\textquotedblright{}.

\begin{cousinweb}
\textbf{\ar{تمساح}} (\emph{timsāḥ}) is listed by Lane under \emph{m-s-ḥ} as the crocodile. It is usually traced to Coptic \emph{msah}, the Egyptian word for the crocodile (EtymDB), with an added \emph{t-}.
\end{cousinweb}

\begin{grammarlens}[title={What you've built}]
One more word for this chapter.
\end{grammarlens}

\subsection*{Your turn}
\begin{itemize}
\raggedright
  \item \emph{Say it:} \emph{timsāḥ}
  \item \emph{Say it:} \emph{timsāḥ}, once more
  \item \emph{Say it:} say \emph{timsāḥ}
  \item \emph{Recall:} say the Arabic for a duck, then the Arabic for a turtle, then say \emph{timsāḥ} again
\end{itemize}

\begin{tcolorbox}[breakable,colback=teal!4,colframe=teal!35!black,title={Before you move on}]
What does \ar{تمساح} mean? (\textquotedblleft{}A crocodile\textquotedblright{}.) Say it once more.
\end{tcolorbox}
\section[\texorpdfstring{\ar{بعوضة} (baʿūḍa)}{بعوضة (baʿūḍa)}]{\ar{بعوضة} (baʿūḍa) — a mosquito}

\label{lesson:AR-C136-bauda}



\begin{quote}
Before the new one: say the Arabic for a turtle, then the Arabic for a crocodile.
\end{quote}

\subsection*{You'll want to know: \ar{بعوضة}}
\textbf{\ar{بعوضة}} — \emph{baʿūḍa} — \textquotedblleft{}a mosquito\textquotedblright{}.

\begin{cousinweb}
\textbf{\ar{بعوضة}} (\emph{baʿūḍa}) is built on the native Arabic root \emph{b-ʿ-ḍ}, \textquotedblleft{}gnats; to cut into parts\textquotedblright{} (Lane). \emph{Baʿūḍa} is one gnat, a mosquito.
\end{cousinweb}

\begin{grammarlens}[title={What you've built}]
One more word for this chapter.
\end{grammarlens}

\subsection*{Your turn}
\begin{itemize}
\raggedright
  \item \emph{Say it:} \emph{baʿūḍa}
  \item \emph{Say it:} \emph{baʿūḍa}, once more
  \item \emph{Say it:} say \emph{baʿūḍa}
  \item \emph{Recall:} say the Arabic for a turtle, then the Arabic for a crocodile, then say \emph{baʿūḍa} again
\end{itemize}

\begin{tcolorbox}[breakable,colback=teal!4,colframe=teal!35!black,title={Before you move on}]
What does \ar{بعوضة} mean? (\textquotedblleft{}A mosquito\textquotedblright{}.) Say it once more.
\end{tcolorbox}
\section[\texorpdfstring{\ar{ذبابة} (dhubāba)}{ذبابة (dhubāba)}]{\ar{ذبابة} (dhubāba) — a fly}

\label{lesson:AR-C136-dhubaba}



\begin{quote}
Before the new one: say the Arabic for a crocodile, then the Arabic for a mosquito.
\end{quote}

\subsection*{You'll want to know: \ar{ذبابة}}
\textbf{\ar{ذبابة}} — \emph{dhubāba} — \textquotedblleft{}a fly\textquotedblright{}.

\begin{cousinweb}
\textbf{\ar{ذبابة}} (\emph{dhubāba}) is built on the native Arabic root \emph{dh-b-b}, \textquotedblleft{}to drive away\textquotedblright{} (Lane), said above all of flies; \emph{dhubāba} is one fly of \emph{dhubāb}. Hebrew \emph{zevuv}, \textquotedblleft{}fly\textquotedblright{}, heard in the name Beelzebub, is its cousin.
\end{cousinweb}

\begin{grammarlens}[title={What you've built}]
One more word for this chapter.
\end{grammarlens}

\subsection*{Your turn}
\begin{itemize}
\raggedright
  \item \emph{Say it:} \emph{dhubāba}
  \item \emph{Say it:} \emph{dhubāba}, once more
  \item \emph{Say it:} say \emph{dhubāba}
  \item \emph{Recall:} say the Arabic for a crocodile, then the Arabic for a mosquito, then say \emph{dhubāba} again
\end{itemize}

\begin{tcolorbox}[breakable,colback=teal!4,colframe=teal!35!black,title={Before you move on}]
What does \ar{ذبابة} mean? (\textquotedblleft{}A fly\textquotedblright{}.) Say it once more.
\end{tcolorbox}
\section[\texorpdfstring{\ar{نحلة} (naḥla)}{نحلة (naḥla)}]{\ar{نحلة} (naḥla) — a bee}

\label{lesson:AR-C136-nahla}



\begin{quote}
Before the new one: say the Arabic for a mosquito, then the Arabic for a fly.
\end{quote}

\subsection*{You'll want to know: \ar{نحلة}}
\textbf{\ar{نحلة}} — \emph{naḥla} — \textquotedblleft{}a bee\textquotedblright{}.

\begin{cousinweb}
\textbf{\ar{نحلة}} (\emph{naḥla}) is one bee of \emph{naḥl}, \textquotedblleft{}bees\textquotedblright{}, an old native noun built on \emph{n-ḥ-l}; Lane's brief article on this root does not list it.
\end{cousinweb}

\begin{grammarlens}[title={What you've built}]
One more word for this chapter.
\end{grammarlens}

\subsection*{Your turn}
\begin{itemize}
\raggedright
  \item \emph{Say it:} \emph{naḥla}
  \item \emph{Say it:} \emph{naḥla}, once more
  \item \emph{Say it:} say \emph{naḥla}
  \item \emph{Recall:} say the Arabic for a mosquito, then the Arabic for a fly, then say \emph{naḥla} again
\end{itemize}

\begin{tcolorbox}[breakable,colback=teal!4,colframe=teal!35!black,title={Before you move on}]
What does \ar{نحلة} mean? (\textquotedblleft{}A bee\textquotedblright{}.) Say it once more.
\end{tcolorbox}
\section[\texorpdfstring{\ar{فراشة} (farāsha)}{فراشة (farāsha)}]{\ar{فراشة} (farāsha) — a butterfly}

\label{lesson:AR-C136-farasha}



\begin{quote}
Before the new one: say the Arabic for a fly, then the Arabic for a bee.
\end{quote}

\subsection*{You'll want to know: \ar{فراشة}}
\textbf{\ar{فراشة}} — \emph{farāsha} — \textquotedblleft{}a butterfly\textquotedblright{}.

\begin{cousinweb}
\textbf{\ar{فراشة}} (\emph{farāsha}) is built on the native Arabic root \emph{f-r-sh}, \textquotedblleft{}to spread out\textquotedblright{} (Lane). Lane's \emph{farāsh} are moths that fly into a lamp, and he adds that modern usage applies it to butterflies.
\end{cousinweb}

\begin{grammarlens}[title={What you've built}]
One more word for this chapter.
\end{grammarlens}

\subsection*{Your turn}
\begin{itemize}
\raggedright
  \item \emph{Say it:} \emph{farāsha}
  \item \emph{Say it:} \emph{farāsha}, once more
  \item \emph{Say it:} say \emph{farāsha}
  \item \emph{Recall:} say the Arabic for a fly, then the Arabic for a bee, then say \emph{farāsha} again
\end{itemize}

\begin{tcolorbox}[breakable,colback=teal!4,colframe=teal!35!black,title={Before you move on}]
What does \ar{فراشة} mean? (\textquotedblleft{}A butterfly\textquotedblright{}.) Say it once more.
\end{tcolorbox}
